\vspace{-2mm}
\section{Introduction}

Image segmentation studies the problem of grouping  pixels.
Different semantics for grouping pixels, \eg, category or instance membership, have led to different types of segmentation tasks, such as panoptic, instance or semantic segmentation.
While these tasks differ only in semantics, current methods develop specialized architectures for each task.
Per-pixel classification architectures
based on Fully Convolutional Networks (FCNs)~\cite{long2015fully} are used for semantic segmentation, while mask classification architectures~\cite{he2017mask,detr} that predict a set of binary masks each associated with a single category,  dominate instance-level segmentation.
Although such \emph{specialized} architectures~\cite{long2015fully,deeplabV3plus,he2017mask,chen2019hybrid} have advanced each individual task, they lack the flexibility to generalize to the other tasks. For example, FCN-based architectures struggle at instance segmentation, leading to the evolution of different architectures for instance segmentation compared to semantic segmentation.
Thus, duplicate research and (hardware) optimization effort is spent on each specialized architecture for every task.

\begin{figure}[t!]
    \centering
    \includegraphics[width=\linewidth]{figures/maskformerv2_teaser.pdf}
    \caption{
        State-of-the-art segmentation architectures are typically specialized for each image segmentation task.
        Although recent work %
        has proposed universal architectures that attempt all tasks and are competitive on semantic and panoptic segmentation, they struggle with segmenting instances. We propose \textbf{\modelname}, which, for the first time, outperforms the best specialized architectures on three studied segmentation tasks on multiple datasets.
        }
    \label{fig:teaser}
    \vspace{-1mm}
\end{figure}

To address this fragmentation, %
recent work~\cite{cheng2021maskformer,zhang2021knet} has attempted to design \emph{universal architectures}, that are capable of addressing all segmentation tasks with the same architecture (\ie, universal image segmentation).
These architectures are typically based on an end-to-end set prediction objective (\eg, DETR~\cite{detr}), and successfully tackle multiple tasks without modifying the architecture, loss, or the training procedure. 
Note, universal architectures are still trained separately for different tasks and datasets, albeit having the same architecture. %
In addition to being flexible, universal architectures have recently shown \sota
results on semantic and panoptic segmentation~\cite{cheng2021maskformer}.
However, recent work  still focuses on advancing specialized architectures~\cite{fapn,strudel2021segmenter,QueryInst}, which raises the question: why haven't universal architectures replaced specialized ones?

Although existing universal architectures are flexible enough to tackle any segmentation task, %
as shown in \figref{fig:teaser}, in practice their performance %
lags behind the best specialized architectures.
For instance, the best reported performance of universal architectures~\cite{cheng2021maskformer,zhang2021knet}, is currently lower ($>9$ AP) than the SOTA specialized architecture for instance segmentation~\cite{chen2019hybrid}.
Beyond the inferior performance, universal architectures are also harder to train. They typically require more advanced hardware and a much longer training schedule. For example, training MaskFormer~\cite{cheng2021maskformer} takes 300 epochs to reach 40.1 AP and it can only fit a single image in a GPU with 32G memory. In contrast, the specialized Swin-HTC++~\cite{chen2019hybrid} obtains better performance in only 72 epochs. Both the performance and training efficiency issues hamper the deployment of universal architectures.


In this work, we propose a universal image segmentation architecture named \modelnamelong (\textbf{\modelname}) that outperforms specialized architectures across different segmentation tasks, while still being easy to train on every task.
We build upon a simple meta architecture~\cite{cheng2021maskformer} consisting of a
backbone feature extractor~\cite{he2016deep,liu2021swin}, a pixel decoder~\cite{lin2016feature} and a Transformer decoder~\cite{vaswani2017attention}.
We propose key improvements that enable better results and efficient training.
First, we use \emph{masked attention} in the Transformer decoder which restricts the attention to localized features centered around predicted segments, which can be either objects or regions depending on the specific semantic for grouping.
Compared to the  cross-attention used in a standard Transformer decoder which attends to all locations in an image, our masked attention leads to faster convergence and improved performance.
Second, we use \emph{multi-scale high-resolution features} which help the model to segment small objects/regions.
Third, we propose \emph{optimization improvements} such as switching the order of self and cross-attention, making query features learnable, and removing dropout; all of which improve performance without additional compute.
Finally, we save $3\x$ training memory without affecting the performance by \emph{calculating mask loss on few randomly sampled points}.
These improvements not only boost the model performance, but also make training significantly easier, making universal architectures more accessible to users with limited compute.

We evaluate \modelname on three image segmentation tasks (panoptic, instance and semantic segmentation) using four popular datasets (COCO~\cite{lin2014coco}, Cityscapes~\cite{Cordts2016Cityscapes}, ADE20K~\cite{zhou2017ade20k} and Mapillary Vistas~\cite{neuhold2017mapillary}).
For the first time, on all these benchmarks, our single architecture performs on par or better than specialized architectures.
\modelname sets the new state-of-the-art of \textbf{57.8 PQ} on COCO panoptic segmentation~\cite{kirillov2017panoptic}, \textbf{50.1 AP} on COCO instance segmentation~\cite{lin2014coco} and \textbf{57.7 mIoU} on ADE20K semantic segmentation~\cite{zhou2017ade20k} using the exact same architecture.


















